% Generated draft: requires mathematical and source review.
\documentclass[11pt]{article}
\usepackage[margin=25mm]{geometry}
\usepackage{fontspec}
\setmainfont{Latin Modern Roman}
\usepackage{amsmath,amssymb,mathtools,unicode-math}
\setmathfont{Latin Modern Math}
\usepackage{enumitem,needspace,xurl,hyperref}
\setlength{\parindent}{0pt}
\setlength{\parskip}{5pt}
\setlength{\emergencystretch}{3em}
\newcommand{\N}{\mathbb{N}}
\newcommand{\Z}{\mathbb{Z}}
\newcommand{\Q}{\mathbb{Q}}
\newcommand{\R}{\mathbb{R}}
\newcommand{\C}{\mathbb{C}}
\newcommand{\F}{\mathbb{F}}
\newcommand{\A}{\mathbb{A}}
\newcommand{\PP}{\mathbb{P}}
\newcommand{\E}{\mathbb{E}}
\newcommand{\eps}{\varepsilon}
\newcommand{\dd}{\,\mathrm{d}}
\begin{document}
\section*{UnsolvedMath ID: 700668}
\textbf{Source problem code:} catalogue entry
\subsection{Definitions and mathematical statement}
Under the standard abelian-group interpretation of the supplied wording, regard \(\Z\) as its additive group and let \(\mathfrak R\), the \emph{Reid class}, be the smallest class of abelian groups, understood up to isomorphism, containing \(\Z\) and satisfying the following closure conditions: for every set \(I\) and every family \((G_i)_{i\in I}\) of members of \(\mathfrak R\),
\[
\bigoplus_{i\in I}G_i\in\mathfrak R
\qquad\text{and}\qquad
\prod_{i\in I}G_i\in\mathfrak R.
\]
Here the direct sum consists of the tuples with finite support, and the direct product consists of all tuples, with componentwise addition. An abelian group \(A\) is a direct summand of \(G\) if there exists an abelian group \(B\) such that \(G\cong A\oplus B\).

The question is whether \(\mathfrak R\) is closed under direct summands; equivalently, whether
\[
\forall A,B\text{ abelian groups},\qquad
A\oplus B\in\mathfrak R\ \Longrightarrow\ A\in\mathfrak R.
\]
The supplied text omits the noun after ``smallest'' and does not explicitly specify the ambient category, the permitted indexing sets, or closure under isomorphism. The formulation above makes the conventional interpretation explicit; these conventions require source confirmation.
\subsection{Short English statement}
Starting with the additive group of integers, form the smallest class closed under arbitrary set-indexed direct sums and direct products. Must every direct summand of a group in this class also belong to the class? The supplied registry labels this question open; that label has not been independently verified.
\subsection{Sources}
\begin{itemize}
\item Local UnsolvedMath registry, permanent ID \(700668\), exact source problem code \texttt{catalogue entry}, title \emph{Kourovka Notebook Problem 10.53}. The supplied record gives status \texttt{open} and proposed year \(1986\).
\item Supplied source attribution: \emph{The Kourovka Notebook (Unsolved Problems in Group Theory)}, 10th issue (1986), Problem 10.53; also identified with the 21st issue edition (2026), arXiv:1401.0300v48. The collection description names E. I. Khukhro and V. D. Mazurov as editors. These bibliographic details are supplied source data and have not been independently checked.
\item Attribution qualification: the statement begins ``(M. Dugas)'', whereas the background and proposer field name S. V. Rychkov. The supplied record does not explain their respective roles.
\end{itemize}
\end{document}
